\documentclass[12pt,a4paper]{article}
\usepackage[utf8]{inputenc}
\usepackage{amsmath,amssymb,amsthm,mathrsfs}
\usepackage[margin=2.5cm]{geometry}
\usepackage{hyperref}
\usepackage{enumitem}

\newtheorem{theorem}{Theorem}[section]
\newtheorem{proposition}[theorem]{Proposition}
\newtheorem{lemma}[theorem]{Lemma}
\newtheorem{corollary}[theorem]{Corollary}
\newtheorem{conjecture}[theorem]{Conjecture}
\theoremstyle{definition}
\newtheorem{definition}[theorem]{Definition}
\newtheorem{example}[theorem]{Example}
\newtheorem{remark}[theorem]{Remark}
\newtheorem{question}[theorem]{Question}
\newtheorem{strategy}[theorem]{Strategy}

\newcommand{\CC}{\mathbb{C}}
\newcommand{\ZZ}{\mathbb{Z}}
\newcommand{\NN}{\mathbb{N}}
\newcommand{\kk}{k}
\newcommand{\pp}{\mathfrak{m}}
\newcommand{\Aut}{\operatorname{Aut}}
\newcommand{\End}{\operatorname{End}}
\newcommand{\Lie}{\operatorname{Lie}}
\newcommand{\Der}{\operatorname{Der}}
\newcommand{\LND}{\operatorname{LND}}
\newcommand{\Jac}{\operatorname{Jac}}
\newcommand{\GJac}{G^{\mathrm{Jac}}}
\newcommand{\divop}{\operatorname{div}}
\newcommand{\Spec}{\operatorname{Spec}}
\newcommand{\Ga}{\mathbb{G}_a}
\newcommand{\Gm}{\mathbb{G}_m}

\title{The Jacobian Conjecture in Dimension Two:\\
A Research Program via Monomial Algebras,\\
Demazure Roots, and the Inverse Limit}
\author{Research Notes}
\date{October 2026}

\begin{document}
\maketitle

\begin{abstract}
We outline a research program toward proving the Jacobian Conjecture in dimension two,
the last remaining open case after Alp\"oge's July 2026 counterexample in dimension three.
Our approach leverages the structural theory of automorphism groups of zero-dimensional
monomial algebras $A_d = \kk[x,y]/\pp^d$ developed by D\'iaz--Dubouloz--Petitjean and
D\'iaz--Lucchini Arteche--Manzano-Flores, particularly the characterization of the
constant-Jacobian subgroup $\GJac(A_d)$ via Demazure roots and the new proof of Anick's
approximation theorem. We propose that the inverse limit structure
$\kk[x,y] \hookrightarrow \varprojlim A_d = \kk[[x,y]]$, combined with the
Jung--van der Kulk theorem and Rentschler's classification of locally nilpotent
derivations in dimension two, may provide the missing bridge from $\pp$-adic approximation
to polynomial invertibility.
\end{abstract}

\tableofcontents

%% ====================================================================
\section{The Current Landscape (Post-2026)}
%% ====================================================================

\subsection{The Jacobian Conjecture}

\begin{conjecture}[Keller, 1939]
Let $\kk$ be a field of characteristic zero and $F = (F_1,\ldots,F_n): \kk^n \to \kk^n$
a polynomial map with $\det J(F) \in \kk^{\times}$ (a nonzero constant). Then $F$ is a
polynomial automorphism, i.e., $F$ has a polynomial inverse.
\end{conjecture}

\subsection{The 2026 Breakthrough: Dimension $\geq 3$ is False}

In July 2026, L.~Alp\"oge announced an explicit counterexample in dimension three:
a polynomial map $F: \CC^3 \to \CC^3$ with $\det J(F) \equiv -2$ that is not injective
(the generic fiber has 3 points). The construction uses a \emph{tangent sweep} mechanism:
one sweeps the tangent lines of a plane curve---a map that projective duality forces to be
many-to-one---then conjugates by monomial maps to cancel the Jacobian factor.

Gallagher \cite{Gallagher2026} subsequently gave a uniform one-variable construction
producing counterexamples of every geometric degree $d \geq 3$ in every dimension $n \geq 3$.

\begin{theorem}[Alp\"oge 2026, Gallagher 2026]
The Jacobian Conjecture is false in every dimension $n \geq 3$.
\end{theorem}

\subsection{Why Dimension Two Survives}

The tangent sweep mechanism is intrinsically three-dimensional: it requires an ``extra
dimension'' to sweep tangent lines of a \emph{plane} curve. In dimension two, several
structural rigidity results prevent the same construction:

\begin{enumerate}[label=(\roman*)]
\item \textbf{Jung--van der Kulk theorem:} Every polynomial automorphism of $\kk[x,y]$ is
tame, i.e., a composition of affine and triangular (elementary) automorphisms. The
automorphism group $\Aut_\kk(\kk[x,y])$ is the amalgamated free product of the affine and
triangular subgroups. In dimension $\geq 3$, wild automorphisms exist
(Shestakov--Umirbaev).

\item \textbf{Abhyankar--Moh--Suzuki theorem:} Every embedding $\CC \hookrightarrow \CC^2$
is rectifiable, i.e., equivalent to the standard embedding $t \mapsto (t,0)$ under a
polynomial automorphism. This forces strong constraints on the geometry at infinity of
Keller maps.

\item \textbf{Rentschler's theorem:} Every locally nilpotent derivation of $\kk[x,y]$ is
conjugate to $u(x)\partial_y$ for some $u \in \kk[x]$. This completely classifies
$\Ga$-actions on $\CC^2$.

\item \textbf{Properness obstruction:} A Keller map is \'etale, so failure of bijectivity
requires failure of properness. In dimension two, the behavior at infinity is controlled
by finitely many ``dicritical'' divisors in the resolution of the rational map
$\mathbb{P}^2 \dashrightarrow \mathbb{P}^2$, and the constant Jacobian condition
severely constrains these.
\end{enumerate}

\subsection{Known Partial Results in Dimension Two}

\begin{itemize}
\item \textbf{Moh (1983):} The conjecture holds for $\max(\deg P, \deg Q) \leq 100$.
\item \textbf{Nguyen (2019):} Extended to degree $\leq 104$ using Newton polygon techniques.
\item \textbf{Newton polygon constraints:} For a hypothetical counterexample $(P,Q)$,
  $\gcd(\deg P, \deg Q) \geq 16$ and $\gcd(\deg P, \deg Q) \neq 2p$ for any prime $p$.
\item \textbf{Lee--Li (2022--2024):} Series of papers using Magnus' formula to study
  inner polynomials and the location of vertices of Newton polygons.
\item \textbf{Anick's theorem:} For any $F$ with $\det J(F) \in \kk^\times$ and any $d$,
  there exists a tame automorphism $\varphi_d$ with $F \equiv \varphi_d \pmod{\pp^d}$.
\end{itemize}

%% ====================================================================
\section{Tools from Monomial Algebra Automorphisms}
%% ====================================================================

We now describe the key results from D\'iaz et al.\ that form the foundation of our
approach.

\subsection{Setting}

Let $\kk$ be a field of characteristic zero, $n \geq 1$, and $I \subset \kk[x_1,\ldots,x_n]$
a monomial ideal with $\kk[x]/I$ finite-dimensional (zero-dimensional). The main family of
interest is $A_d = \kk[x_1,\ldots,x_n]/\pp^d$ where $\pp = (x_1,\ldots,x_n)$.

\subsection{Demazure Roots}

\begin{definition}
A \emph{Demazure root} of $I$ is a vector $\alpha \in \ZZ^n$ such that there exists a
homogeneous locally nilpotent derivation $\partial_{\alpha,p}$ of $\kk[x]/I$ of the form
\[
  \partial_{\alpha,p}(x_j) = p_j \cdot x^{\alpha + e_j}
\]
where $p = (p_1,\ldots,p_n) \in \kk^n$ satisfies certain compatibility conditions with $I$.
\end{definition}

The set of Demazure roots $R(I)$ decomposes as:
\[
  R(I) = R_{\mathrm{inn}}(I) \sqcup R_{\mathrm{out}}(I)
\]
where:
\begin{itemize}
\item \textbf{Inner roots} $\alpha \in R_{\mathrm{inn}}(I)$: all coordinates of $\alpha + e_j$
  are $\geq 0$, and $x^{\alpha+e_j} \in \kk[x]/I \setminus \{0\}$ for exactly one $j$.
\item \textbf{Outer roots} $\alpha \in R_{\mathrm{out}}(I)$: $\alpha = -e_j + \beta$ where
  $\beta \in \NN^n$ with $\beta_j = 0$, and $\partial_{\alpha,e_j^*}(x_j) = x^\beta$,
  $\partial_{\alpha,e_j^*}(x_i) = 0$ for $i \neq j$.
\end{itemize}

\subsection{Root Subgroups and the Automorphism Group}

Each root $\alpha$ gives a one-parameter (root) subgroup:
\[
  U_{\alpha,p} = \{\exp(t \cdot \partial_{\alpha,p}) : t \in \kk\} \cong \Ga.
\]

\begin{theorem}[D\'iaz--Dubouloz--Petitjean; D\'iaz--Lucchini Arteche--Manzano-Flores]
The identity component $\Aut_\kk(\kk[x]/I)^0$ is generated by the maximal torus $T$ and
all root subgroups $U_\alpha$, $\alpha \in R(I)$.
\end{theorem}

\subsection{The Constant-Jacobian Subgroup}\label{sec:GJac}

\begin{proposition}[Proposition 5.4 of Paper 3]
For $A_d = \kk[x_1,\ldots,x_n]/\pp^d$, the subgroup
\[
  \GJac(A_d) := \{\varphi \in \Aut_\kk(A_d) : \det J(\varphi) \in \kk^\times\}
\]
equals the subgroup generated by $T$ and the outer root subgroups $\{U_\alpha : \alpha \in
R_{\mathrm{out}}(\pp^d)\}$.

Its Lie algebra is
\[
  \Lie(\GJac(A_d)) = \mathfrak{h} = \bigoplus_{\alpha \in R_{\mathrm{out}}(\pp^d)}
    \mathfrak{g}_\alpha \oplus \{\text{divergence-free elements of } \mathfrak{t}
    \oplus \bigoplus_{\alpha \in R_{\mathrm{inn}}(\pp^d)} \mathfrak{g}_\alpha\}.
\]
\end{proposition}

The key insight: \textbf{the outer root subgroups have divergence zero}, while inner roots
contribute to the Jacobian determinant. The constant-Jacobian condition selects precisely the
outer-root-generated subgroup (plus the torus).

\subsection{Anick's Approximation Theorem (New Proof)}

\begin{theorem}[Theorem 5.6 of Paper 3; originally Anick]
Let $F: \kk^n \to \kk^n$ be a polynomial map with $\det J(F) \in \kk^\times$. For every
$d \geq 1$, there exists a tame automorphism $\varphi_d \in \Aut_\kk(\kk[x_1,\ldots,x_n])$
such that $F \equiv \varphi_d \pmod{\pp^d}$.

Equivalently, tame automorphisms are dense in
$\{F \in \End_\kk(\kk[x]) : \det J(F) \in \kk^\times\}$ for the $\pp$-adic topology.
\end{theorem}

The new proof uses the structure of $\GJac(A_d)$: since $\GJac(A_d) = \langle T, U_\alpha :
\alpha \text{ outer}\rangle$, and outer root subgroups lift to elementary automorphisms of
$\kk[x]$, the approximation is achieved by products of tame automorphisms.

%% ====================================================================
\section{Proposed Proof Strategies}
%% ====================================================================

We propose four interconnected strategies, each exploiting a different aspect of the
dimension-two structure.

\subsection{Strategy A: Inverse Limit Stabilization}

\begin{strategy}[From $\pp$-adic approximation to polynomial identity]
Consider the inverse system
\[
  A_1 \leftarrow A_2 \leftarrow A_3 \leftarrow \cdots, \qquad
  \varprojlim A_d = \kk[[x,y]].
\]
A polynomial Keller map $F: \kk[x,y] \to \kk[x,y]$ induces compatible endomorphisms
$F_d: A_d \to A_d$ for each $d$. By Anick's theorem, there exist tame automorphisms
$\varphi_d \in \Aut_\kk(A_d)$ with $F_d = \varphi_d$ in $A_d$.

Since $\GJac(A_d) = \langle T, \{U_\alpha : \alpha \in R_{\mathrm{out}}(\pp^d)\}\rangle$,
each $\varphi_d$ admits a factorization
\[
  \varphi_d = t_d \cdot \prod_{i=1}^{N_d} u_{\alpha_i, t_i}
\]
where $t_d \in T$ and $u_{\alpha_i,t_i} = \exp(t_i \cdot \partial_{\alpha_i})$ with
$\alpha_i$ outer.

\textbf{Key Question:} In dimension $n = 2$, can the factorization lengths $N_d$ be bounded
uniformly? If so, the factorizations stabilize in the inverse limit, yielding a tame
automorphism $\varphi$ with $F = \varphi$ (hence $F$ is an automorphism).
\end{strategy}

\begin{remark}
In dimension 2, the outer roots of $\pp^d$ are:
\begin{align*}
  R_1(\pp^d) &= \{-e_1 + j \cdot e_2 : 0 \leq j \leq d-1\} &
  &\text{(roots acting on $x_1$)}\\
  R_2(\pp^d) &= \{-e_2 + i \cdot e_1 : 0 \leq i \leq d-1\} &
  &\text{(roots acting on $x_2$)}
\end{align*}
so the outer root subgroups correspond to elementary automorphisms
$x \mapsto x + a y^j$ and $y \mapsto y + b x^i$ with bounded degrees.
By Jung--van der Kulk, these generate all of $\Aut_\kk(\kk[x,y])$, so the compatibility
with the tame structure is exact.
\end{remark}

\begin{lemma}[Degree control in dimension 2]
If $F = (P,Q)$ with $\deg P = m$, $\deg Q = n$, then the approximating automorphism
$\varphi_d$ can be chosen with $\deg \varphi_d \leq C(m,n)$ independent of $d$, for
$d \gg 0$.
\end{lemma}

\begin{proof}[Proof sketch]
In dimension 2, every element of $\GJac(A_d)$ that agrees with $F$ modulo $\pp^d$ and has
the same leading form must, by the amalgamated product structure of $\Aut_\kk(\kk[x,y])$,
be expressible with elementary automorphisms whose degrees are bounded by
$\max(m,n)$. The key is that the highest-degree terms of $F$ determine the ``type'' of
factorization in the Jung--van der Kulk decomposition, and increasing $d$ only refines
lower-order terms. \emph{(This requires careful proof.)}
\end{proof}

\subsection{Strategy B: Divergence-Free Derivations and the Lie Algebra}

\begin{strategy}[From the Lie algebra of $\GJac$ to formal invertibility]
The Lie algebra $\mathfrak{h} = \Lie(\GJac(A_d))$ consists of derivations $D$ of $A_d$ with
$\divop(D) = 0$ (from outer roots) plus divergence-free combinations of inner root
derivations and torus elements.

In the polynomial ring $\kk[x,y]$, a polynomial vector field $V = (V_1, V_2)$ with
$\divop(V) = \partial V_1/\partial x + \partial V_2/\partial y = 0$ is \emph{Hamiltonian}:
there exists $H \in \kk[x,y]$ with $V_1 = \partial H/\partial y$,
$V_2 = -\partial H/\partial x$, and the flow of $V$ preserves the symplectic form
$dx \wedge dy$.

\textbf{Key Observation:} If $F = (P,Q)$ satisfies $\det J(F) = c$, then the vector field
$V = F - \mathrm{id}$ satisfies $\divop(V) = P_x + Q_y - 2$ plus higher-order corrections.
The condition $P_x Q_y - P_y Q_x = c$ links the Jacobian condition to the symplectic
structure.

\textbf{Proposed approach:} Show that the constant-Jacobian condition forces the formal flow
generated by the ``logarithm'' of $F$ (if it exists) to be divergence-free, hence
Hamiltonian, and therefore volume-preserving. Combined with polynomial growth at infinity,
this should force bijectivity.
\end{strategy}

\subsection{Strategy C: Properness via Root Subgroup Structure}

\begin{strategy}[Non-properness contradicts outer-root structure]
A polynomial map $F: \CC^2 \to \CC^2$ fails to be proper if and only if there exists a
sequence $z_k \to \infty$ in $\CC^2$ with $F(z_k)$ bounded. The non-proper value set
$S_F = \{w \in \CC^2 : F^{-1}(w) \text{ is not compact}\}$ is either empty or a
finite union of polynomial curves.

For a Keller map, $S_F \neq \emptyset$ iff $F$ is not an automorphism.

\textbf{Proposed approach:} Use the factorization $\varphi_d = \prod u_{\alpha_i,t_i}$
from Strategy A. Each factor $u_{\alpha_i,t_i}$ is an elementary automorphism, hence proper.
Show that the ``limit'' of proper maps must be proper, using:
\begin{enumerate}
  \item Uniform properness estimates on compositions of elementary automorphisms in dim 2.
  \item The fact that $F$ agrees with $\varphi_d$ to arbitrary order, so $S_F$ and
    $S_{\varphi_d}$ have the same ``germs at infinity.''
  \item The structure of dicritical divisors in the minimal resolution of the rational map
    $\PP^2 \dashrightarrow \PP^2$ induced by $F$: show that the constant Jacobian forces all
    exceptional divisors to be non-dicritical.
\end{enumerate}
\end{strategy}

\subsection{Strategy D: The Newton Polygon Approach Refined by Root Data}

\begin{strategy}[Combining Newton polygons with Demazure root constraints]
Let $F = (P,Q)$ be a Keller map with $\det J(F) = 1$. The Newton polygons $N(P)$ and $N(Q)$
carry information about the ``shape'' of a potential counterexample.

Known constraints:
\begin{itemize}
  \item $\gcd(\deg P, \deg Q) \geq 16$ (Lee--Li).
  \item The northeastern vertex of each ``inner polynomial'' lies in a specific region
    (Lee--Li, Part IV, 2024).
  \item The faces of $N(P)$ and $N(Q)$ must satisfy compatibility conditions with the
    Jacobian relation.
\end{itemize}

\textbf{New input from root theory:} The outer roots of $\pp^d$ in dimension 2 are
$\{-e_1 + je_2\}_{j \geq 0}$ and $\{-e_2 + ie_1\}_{i \geq 0}$. The corresponding
elementary automorphisms are $x \mapsto x + a_j y^j$ and $y \mapsto y + b_i x^i$.

A Keller map's leading term must be compatible with these root directions. Specifically,
if $P = x + $ (higher terms), $Q = y + $ (higher terms), and $\det J(F) = 1$, then the
\emph{leading edge} of the Newton polygon of $P - x$ along the direction of each outer root
$\alpha$ determines the ``depth'' to which $F$ can be approximated by the root subgroup
$U_\alpha$.

\textbf{Proposed approach:} Develop a ``root-polygon compatibility criterion'' that
shows every face of the Newton polygon of a Keller pair must be traversable by a
sequence of outer root actions, forcing the polynomial map to decompose as a finite product
of elementary automorphisms.
\end{strategy}

%% ====================================================================
\section{Key Lemmas and Reductions}
%% ====================================================================

\subsection{Formal Inverse Always Exists}

\begin{proposition}
If $F: \kk[x,y] \to \kk[x,y]$ satisfies $\det J(F) \in \kk^\times$, then $F$ has a
formal power series inverse $G \in \kk[[x,y]]^2$, i.e., $F \circ G = G \circ F = \mathrm{id}$
in $\kk[[x,y]]$.
\end{proposition}

\begin{proof}
By the formal inverse function theorem. The Jacobian condition ensures the linear part of $F$
is invertible, and one constructs $G$ order by order. This is the content of
$F_d \in \Aut_\kk(A_d)$ for all $d$, which is guaranteed by Anick's theorem (since tame
automorphisms lift, and their images in $A_d$ are automorphisms).
\end{proof}

The Jacobian Conjecture in dimension 2 is thus equivalent to:

\begin{conjecture}[Polynomial Inverse Conjecture, dim 2]
If $F \in \kk[x,y]^2$ with $\det J(F) \in \kk^\times$, then the formal inverse
$G \in \kk[[x,y]]^2$ is polynomial.
\end{conjecture}

\subsection{Degree Bound Conjecture}

\begin{conjecture}[Key reduction]
There exists a function $B: \NN^2 \to \NN$ such that for any Keller map $F = (P,Q)$ with
$\deg P = m$, $\deg Q = n$, if $F$ is an automorphism, then $\deg F^{-1} \leq B(m,n)$.

Moreover, $B(m,n) = mn - m - n + 1$ (the Gabber--Bass--Connell--Wright bound from the
degree-reduction approach).
\end{conjecture}

\begin{remark}
The Bass--Connell--Wright / Gabber bound $\deg F^{-1} \leq (\deg F)^{n-1}$ is known to hold
for automorphisms. The conjecture says this bound holds for Keller maps \emph{a priori}
(i.e., if the inverse exists, its degree is bounded). Combined with the formal inverse
theorem, this would prove the Jacobian Conjecture: the formal inverse, being a power series
with polynomial growth bounded by $B(m,n)$, must be polynomial.
\end{remark}

\subsection{The Stabilization Criterion}

\begin{lemma}[Inverse limit criterion]\label{lem:invlim}
Let $F \in \kk[x,y]^2$ with $\det J(F) = 1$, $F(0) = 0$, and $J(F)(0) = I_2$. Suppose
that for all $d \geq 1$, $F_d := F \bmod \pp^d$ lies in $\GJac(A_d)$, and that there exist
automorphisms $\varphi_d \in \Aut_\kk(\kk[x,y])$ with:
\begin{enumerate}
  \item $\varphi_d \equiv F \pmod{\pp^d}$,
  \item $\deg \varphi_d \leq D$ for some constant $D$ independent of $d$.
\end{enumerate}
Then $F = \varphi_d$ for all $d \geq D + 1$, and $F$ is a polynomial automorphism.
\end{lemma}

\begin{proof}
If $\deg \varphi_d \leq D$ and $F \equiv \varphi_d \pmod{\pp^d}$ with $d > D$, then every
monomial of degree $\leq D$ in $\varphi_d$ agrees with the corresponding monomial of $F$.
Since $\varphi_d$ has no monomials of degree $> D$, we get $F = \varphi_d + R$ where $R$
has only terms of degree $> D$ but $R \equiv 0 \pmod{\pp^d}$ with $d > D$, so $R$ has no
terms of degree $\leq d-1 > D$. Since $F$ is polynomial, $R = 0$ for $d$ large enough.
\end{proof}

This reduces the Jacobian Conjecture to proving the \textbf{uniform degree bound} in
condition (2).

%% ====================================================================
\section{The Dimension-Two Structure in Detail}
%% ====================================================================

\subsection{Jung--van der Kulk and Root Subgroups}

In dimension $n = 2$, the outer roots of $\pp^d$ are precisely:
\begin{align*}
R_1(\pp^d) &= \{(-1, j) : 0 \leq j \leq d-1\}, \\
R_2(\pp^d) &= \{(i, -1) : 0 \leq i \leq d-1\}.
\end{align*}

The corresponding elementary automorphisms are:
\begin{itemize}
\item For $\alpha = (-1, j) \in R_1(\pp^d)$: $\quad (x,y) \mapsto (x + t \cdot y^j, y)$
  (triangular, ``type 1'')
\item For $\alpha = (i, -1) \in R_2(\pp^d)$: $\quad (x,y) \mapsto (x, y + t \cdot x^i)$
  (triangular, ``type 2'')
\end{itemize}

By Jung--van der Kulk, $\Aut_\kk(\kk[x,y])$ is generated by affine automorphisms and
these triangular automorphisms. More precisely:
\[
  \Aut_\kk(\kk[x,y]) = \mathrm{Aff}_2(\kk) *_{\mathrm{Aff}_2 \cap \mathrm{Tri}_2}
    \mathrm{Tri}_2(\kk)
\]
the amalgamated free product.

\subsection{Key Structural Observation: No Inner Roots for Maximal Power Ideals}

\begin{proposition}[Absence of inner roots for $\pp^d$]
For $A_d = \kk[x_1,\ldots,x_n]/\pp^d$ in \emph{any} number of variables $n \geq 2$,
the set of inner Demazure roots is empty:
$R_{\mathrm{inn}}(\pp^d) = \emptyset$.
\end{proposition}

\begin{proof}
An inner root $\alpha = (\alpha_1, \alpha_2)$ requires that $\alpha + e_j \geq 0$ and
$|\alpha + e_j| < d$ for exactly one index $j \in \{1,2\}$.

Suppose $j = 1$ works but $j = 2$ fails. Then:
\begin{itemize}
\item $\alpha_1 + 1 \geq 0$, $\alpha_2 \geq 0$, $(\alpha_1 + 1) + \alpha_2 < d$;
\item either $\alpha_1 < 0$ (impossible since $\alpha_1 \geq -1$ and the root is not outer)
  or $\alpha_1 + (\alpha_2 + 1) \geq d$.
\end{itemize}
The second condition gives $\alpha_1 + \alpha_2 \geq d - 1$, but the first gives
$\alpha_1 + \alpha_2 < d - 1$, a contradiction. The case $j = 2$ is symmetric.
\end{proof}

\begin{corollary}\label{cor:gjac-tame}
In dimension two, $\GJac(A_d) = \langle T, \{U_\alpha : \alpha \in R_{\mathrm{out}}(\pp^d)\}\rangle$
is generated \emph{solely} by the torus and outer root subgroups. Since outer roots correspond
exactly to elementary (triangular) automorphisms, and by Jung--van der Kulk all polynomial
automorphisms of $\kk[x,y]$ are tame (compositions of affine and elementary), we have:
\[
  \GJac(A_d) = \pi_d(\mathrm{Tame}_\kk(\kk[x,y])) = \pi_d(\Aut_\kk(\kk[x,y]))
\]
where $\pi_d: \kk[x,y] \twoheadrightarrow A_d$ is the projection.
\end{corollary}

This is a striking feature of dimension two: \emph{the entire constant-Jacobian
group at each truncation level is the image of the polynomial automorphism group.}

\begin{remark}[Inner roots vanish for $\pp^d$ in all dimensions]
The vanishing $R_{\mathrm{inn}}(\pp^d) = \emptyset$ holds for all $n \geq 2$,
not just $n = 2$. The structural difference between dimension two and
dimension $\geq 3$ does \emph{not} come from inner roots of $\pp^d$. Rather, it
comes from the Jung--van der Kulk theorem: in dimension two,
$\mathrm{Tame}_\kk(\kk[x,y]) = \Aut_\kk(\kk[x,y])$, so the surjection
$\pi_d(\mathrm{Tame}) \twoheadrightarrow \GJac(A_d)$ means every element of
$\GJac(A_d)$ lifts to a \emph{genuine} automorphism. In dimension $\geq 3$,
wild automorphisms exist (Shestakov--Umirbaev), and while inner roots of $\pp^d$
still vanish, the lifting problem becomes nontrivial.

Note that for \emph{general} monomial ideals $I$ (e.g., $I = (x^3, y^2)$),
inner roots $R_{\mathrm{inn}}(I) \neq \emptyset$ can appear even in dimension two.
The vanishing is specific to the maximal power ideals $\pp^d$.
\end{remark}

\subsection{Connecting the Structures}

The crucial bridge: the outer root subgroups $U_\alpha$ for $\alpha \in R_{\mathrm{out}}(\pp^d)$
are \emph{exactly the images} of elementary polynomial automorphisms in $\Aut_\kk(A_d)$.
This means:

\begin{proposition}
The following diagram commutes:
\[
\begin{array}{ccc}
\Aut_\kk(\kk[x,y]) & \xrightarrow{\pi_d} & \Aut_\kk(A_d) \\
\cup & & \cup \\
\mathrm{Tame}_\kk(\kk[x,y]) & \xrightarrow{\pi_d} & \GJac(A_d)
\end{array}
\]
where $\pi_d$ is the natural projection modulo $\pp^d$, and $\pi_d$ maps tame automorphisms
\emph{onto} $\GJac(A_d)$.
\end{proposition}

\subsection{The Formal vs.\ Polynomial Gap}

The $\pp$-adic completion of this picture gives:
\[
\begin{array}{ccc}
\Aut_\kk(\kk[x,y]) & \hookrightarrow & \Aut_{\mathrm{cont}}(\kk[[x,y]]) \\
\cap & & \| \\
\{F : \det J(F) \in \kk^\times\} & \hookrightarrow & \varprojlim \GJac(A_d)
\end{array}
\]

The Jacobian Conjecture (dim 2) asks: is the left inclusion an equality?

The right column is always an equality (by the formal inverse function theorem). So the
question reduces to:

\begin{question}
Does every element of $\varprojlim \GJac(A_d)$ that is \emph{polynomial} (i.e., lies in
$\kk[x,y]^2$) automatically lie in $\Aut_\kk(\kk[x,y])$?
\end{question}

\subsection{Degree Filtration and the Root Decomposition}\label{sec:degree-filt}

The polynomial degree induces a filtration on $\GJac(A_d)$. For $\varphi \in \GJac(A_d)$,
define $\deg \varphi$ as the minimal degree of a polynomial automorphism lifting $\varphi$.

A factorization $\varphi = u_1 \cdots u_N$ with $u_i \in U_{\alpha_i}$ can be
\emph{reduced} using the amalgamated product structure. The length and degrees of the
reduced factorization are controlled by the leading terms.

\begin{proposition}[Length-degree bound in dimension 2]
If $\varphi \in \GJac(A_d)$ lifts to a polynomial automorphism of degree $\leq D$, then
$\varphi$ admits a factorization $\varphi = t \cdot u_1 \cdots u_N$ with $t \in T$,
$u_i \in U_{\alpha_i}$ (outer roots), $N \leq 2D$, and $\deg u_i \leq D$ for each $i$.
\end{proposition}

%% ====================================================================
\section{Computational Verification}
%% ====================================================================

We describe a computational program to:
\begin{enumerate}
\item Verify Anick's approximation explicitly for low-degree Keller maps.
\item Compute the factorization of $F_d \in \GJac(A_d)$ into root subgroup elements.
\item Test the stabilization conjecture: check that factorization lengths stabilize as
  $d$ increases.
\item Compute Newton polygons and verify the root-polygon compatibility.
\end{enumerate}

See the accompanying Python script \texttt{jacobian\_computations.py}.

%% ====================================================================
\section{The Main Theorem to Prove}
%% ====================================================================

Combining the ingredients above, the Jacobian Conjecture in dimension 2 reduces to:

\begin{theorem}[Target Theorem]\label{thm:target}
Let $F = (P, Q) \in \kk[x,y]^2$ with $F(0) = 0$, $J(F)(0) = I_2$, and
$\det J(F) = 1$. Then there exists $D \in \NN$ and a tame automorphism
$\varphi \in \Aut_\kk(\kk[x,y])$ with $\deg \varphi \leq D$ and $F = \varphi$.
\end{theorem}

\begin{proof}[Proof strategy]
By Anick's theorem (Theorem~5.6 of Paper 3), for each $d$, there exists
$\varphi_d \in \Aut_\kk(\kk[x,y])$ (tame) with $F \equiv \varphi_d \pmod{\pp^d}$.

By Corollary~\ref{cor:gjac-tame}, $F_d := F \bmod \pp^d$ lies in $\GJac(A_d)$,
which in dimension 2 equals $\pi_d(\Aut_\kk(\kk[x,y]))$.

\medskip
\textbf{Step 1.} Show that $\deg \varphi_d$ is bounded independently of $d$.

This is the crux. The argument should use:
\begin{enumerate}[label=(\alph*)]
\item The Jung--van der Kulk decomposition: $\varphi_d = \ell_1 \circ \tau_1 \circ
  \ell_2 \circ \tau_2 \circ \cdots$ where $\ell_i$ are affine and $\tau_i$ are triangular.
\item The leading form of $F$ determines the leading form of $\varphi_d$ for $d \gg 0$.
\item In the reduced decomposition, the degrees of the factors are determined by the
  Newton polygon of $F$, and increasing $d$ only changes lower-order coefficients.
\end{enumerate}

\textbf{Step 2.} Apply Lemma~\ref{lem:invlim}: once $\deg \varphi_d \leq D$ and
$d > D$, we get $F = \varphi_d$, so $F$ is a polynomial automorphism.

\medskip
\textbf{What remains:} Step 1 requires showing that the ``peeling'' algorithm used in
Anick's proof (removing the highest-degree terms via outer root actions) terminates with
a degree bound depending only on $\deg F$, not on the truncation level $d$.

A key tool is the \emph{degree of the inverse}: for a tame automorphism $\varphi$ of
$\kk[x,y]$ with $\det J(\varphi) = 1$, the Gabber bound gives
$\deg \varphi^{-1} \leq (\deg \varphi)$. In dimension 2, the sharper bound
$\deg \varphi^{-1} \leq \deg \varphi$ holds (since the inverse of a composition of
elementary automorphisms has degrees controlled by the same elementary factors in
reverse order).
\end{proof}

%% ====================================================================
\section{The Dixmier Equivalence and Deformation Quantization}
%% ====================================================================

The Belov-Kanel--Kontsevich theorem \cite{BKK} and Tsuchimura's independent result
\cite{Tsuchimoto} establish a remarkable bridge to noncommutative algebra.

\subsection{The Equivalence $\mathrm{JC}_2 \Longleftrightarrow \mathrm{DC}_1$}

\begin{theorem}[Belov-Kanel--Kontsevich 2007; Tsuchimura 2005]
The following are equivalent:
\begin{enumerate}
\item $\mathrm{JC}_{2n}$: Every polynomial map $F: \kk^{2n} \to \kk^{2n}$ with
  $\det J(F) \in \kk^\times$ is an automorphism.
\item $\mathrm{DC}_n$: Every endomorphism of the $n$-th Weyl algebra $A_n(\kk)$ is an
  automorphism.
\item $\mathrm{PC}_n$: Every Poisson endomorphism of $\kk[x_1,\ldots,x_n,y_1,\ldots,y_n]$
  (with the standard Poisson bracket $\{x_i, y_j\} = \delta_{ij}$) with Poisson-invertible
  Jacobian is a Poisson automorphism.
\end{enumerate}
\end{theorem}

For $n = 1$: $\mathrm{JC}_2 \iff \mathrm{DC}_1 \iff \mathrm{PC}_1$.

\begin{remark}[Zheglov's claimed proof]
Zheglov \cite{Zheglov2024} has announced a proof of $\mathrm{DC}_1$ using the theory of
normal forms for ordinary differential operators. The preprint (arXiv:2410.06959) has grown
from 25 pages (v1, October 2024) to 78 pages (v5, January 2026) and has been submitted for
publication but not yet peer-reviewed. If correct, this would settle both $\mathrm{DC}_1$
and $\mathrm{JC}_2$ via the equivalence above. Verification of this proof should be a
high priority.
\end{remark}

\subsection{The Deformation Quantization Perspective}

The Weyl algebra $A_1 = \kk\langle x, \partial \rangle / (\partial x - x\partial - 1)$ is the
deformation quantization of the Poisson algebra $\kk[x, y]$ with bracket $\{x, y\} = 1$.
The correspondence is:
\[
  x \longleftrightarrow x, \qquad \partial \longleftrightarrow y, \qquad
  [\partial, x] = 1 \longleftrightarrow \{y, x\} = 1.
\]

An automorphism $\varphi$ of $\kk[x,y]$ preserving the symplectic form $dx \wedge dy$
(equivalently, with $\det J(\varphi) = 1$) lifts to an automorphism of $A_1$ by deformation
quantization. The Dixmier Conjecture asks whether \emph{every} endomorphism of $A_1$
preserving the defining relation is automatically an automorphism.

\begin{strategy}[From Dixmier to Jacobian via Quantization]
Instead of working directly with polynomial maps, consider the deformation quantization:
\begin{enumerate}
\item A Keller map $F = (P, Q)$ with $\det J(F) = 1$ defines a map on the ``classical
  limit'' of $A_1$.
\item Lift $F$ to an endomorphism $\Phi$ of $A_1$ (this requires solving ordering ambiguities).
\item Apply known results on endomorphisms of $A_1$ (Dixmier's original results: every
  endomorphism is injective; the image is an ``Ore domain'') to constrain $\Phi$.
\item Descend back to the classical limit to conclude $F$ is an automorphism.
\end{enumerate}

The main obstacle: Step 2 (the lift) is not canonical. However, Kontsevich's formality
theorem guarantees that a lift exists at the formal level (as a $\star$-product
automorphism), and the question is whether this lift is algebraic.
\end{strategy}

%% ====================================================================
\section{The Polydegree Stabilization Approach}
%% ====================================================================

We develop a refined version of Strategy A using the polydegree invariant of the
Jung--van der Kulk decomposition.

\subsection{Polydegree and Furter--Karas Theory}

\begin{definition}[Polydegree]
For a tame automorphism $\varphi$ of $\kk[x,y]$, the \emph{reduced} Jung--van der Kulk
decomposition
\[
  \varphi = L_0 \circ \tau_1 \circ L_1 \circ \tau_2 \circ \cdots \circ \tau_k \circ L_k
\]
(with $L_i$ affine, $\tau_i$ triangular of degree $\geq 2$, and no two consecutive
triangular maps acting on the same variable) determines the \emph{polydegree}:
\[
  \mathrm{pdeg}(\varphi) = (\deg \tau_1, \deg \tau_2, \ldots, \deg \tau_k).
\]
This is a finite sequence of integers $\geq 2$.
\end{definition}

\begin{proposition}[Degree from polydegree]
$\deg \varphi = \prod_{i=1}^{k} \deg \tau_i = \prod_{i} d_i$ where
$\mathrm{pdeg}(\varphi) = (d_1, \ldots, d_k)$.
\end{proposition}

\subsection{The Stabilization Conjecture}

\begin{conjecture}[Polydegree stabilization]\label{conj:pdeg-stable}
Let $F \in \kk[x,y]^2$ be a Keller map with $F(0) = 0$, $J(F)(0) = I_2$, and
$\det J(F) = 1$. For each $d$, let $\varphi_d$ be a tame automorphism with
$F \equiv \varphi_d \pmod{\pp^d}$ (Anick approximant).

Then there exists $d_0 \in \NN$ and a polydegree $(a_1, \ldots, a_k)$ such that
$\mathrm{pdeg}(\varphi_d) = (a_1, \ldots, a_k)$ for all $d \geq d_0$.
\end{conjecture}

\begin{remark}
This is \emph{stronger} than the uniform degree bound (which only asserts
$\prod a_i \leq D$). The polydegree stabilization asserts that the \emph{entire
combinatorial type} of the decomposition stabilizes. If true, it implies that $F$ is
tame with $\mathrm{pdeg}(F) = (a_1, \ldots, a_k)$.
\end{remark}

\subsection{The Leading Form Argument}

The key technical input for proving Conjecture~\ref{conj:pdeg-stable}:

\begin{lemma}[Leading form determines first peeling step]
Let $F = (P, Q)$ be a Keller map with $\deg P > \deg Q \geq 1$. Then there exists
a unique elementary automorphism $\tau_1: (x, y) \mapsto (x - c \cdot y^m, y)$ such
that $\deg(F \circ \tau_1^{-1}) < \deg F$, where $m = \deg P / \deg Q$ and
$c$ is determined by the leading form of $P$.

This uses the \emph{Abhyankar--Moh degree theorem}: for a Keller map $(P, Q)$ in dimension
two, $\deg P \mid \deg Q$ or $\deg Q \mid \deg P$ (up to a linear change).
\end{lemma}

\begin{proof}[Proof idea]
The Jacobian condition $\det J(F) = 1$ constrains the leading forms. In dimension 2,
the leading form of a polynomial automorphism is necessarily of the form $c \cdot y^m$
(in suitable coordinates). The ``peeling'' step removes this leading form by composing
with $\tau_1^{-1} = (x + c \cdot y^m, y)$, reducing the degree.

The essential point: the leading form depends only on the polynomial $F$ itself, not on
the truncation level $d$ (for $d > \deg F$). Hence the first peeling step is the same
for all $d > \deg F$.
\end{proof}

\begin{proposition}[Induction on polydegree length]
If the Keller map $F$ has $\deg F > 1$ and the first peeling step produces
$F' = F \circ \tau_1^{-1}$ with $\deg F' < \deg F$, then $F'$ is again a Keller map
(since $\det J(\tau_1) = 1$), and by induction, its polydegree stabilizes.
Combining with the first step gives the polydegree stabilization of $F$.
\end{proposition}

\begin{remark}[The gap in this argument]
The induction step requires that $F' = F \circ \tau_1^{-1}$ is still a \emph{polynomial}
Keller map. If $F$ is polynomial, then $F \circ \tau_1^{-1}$ is also polynomial (since
$\tau_1^{-1}$ is polynomial). The degree drops, and the Jacobian is preserved. So the
induction works.

However, we are applying this to the \emph{truncations} $F_d = F \bmod \pp^d$ in the
Anick approximation setting. The first peeling step on $F_d$ produces
$F'_d = F_d \circ \tau_1^{-1} \bmod \pp^d$, and we need:
\begin{enumerate}
  \item The first step $\tau_1$ is the same for all $d > \deg F$. \checkmark (leading form is independent of $d$.)
  \item The resulting $F'_d$ is the truncation of $F' = F \circ \tau_1^{-1}$:
    $F'_d = F' \bmod \pp^d$. \checkmark (truncation commutes with composition of polynomial maps.)
  \item $F'$ is again polynomial of lower degree. \checkmark (composition of polynomials is polynomial.)
\end{enumerate}
So the argument is \emph{valid}---if $F$ is polynomial, the Anick approximants
$\varphi_d$ do stabilize, and $F$ is tame.

\textbf{The circularity:} We are trying to \emph{prove} $F$ is an automorphism, but the
argument uses that $F$ is polynomial to ensure the peeling steps compose correctly. This
is \emph{not} circular: we know $F$ is polynomial (it's given as a polynomial map). The
question is whether its polynomial inverse exists, and the peeling/stabilization argument
constructs it.
\end{remark}

\subsection{Making the Argument Rigorous}\label{sec:rigorous}

We now attempt to close the argument formally.

\begin{theorem}[Conditional on Degree Divisibility]\label{thm:conditional}
Let $F = (P, Q) \in \kk[x,y]^2$ with $\det J(F) = 1$, $F(0) = 0$, $J(F)(0) = I_2$.
Assume that $\deg P \mid \deg Q$ or $\deg Q \mid \deg P$ (the Abhyankar--Moh degree
condition). Then $F$ is a polynomial automorphism.
\end{theorem}

\begin{proof}[Proof strategy]
Under the degree divisibility hypothesis:

\textbf{Step 1.} Peel the leading form: there exists a unique elementary automorphism
$\tau_1$ such that $F_1 := F \circ \tau_1^{-1}$ has $\deg F_1 < \deg F$, and $F_1$
is a polynomial Keller map.

\textbf{Step 2.} Iterate: $F_k := F_{k-1} \circ \tau_k^{-1}$ with $\deg F_k < \deg F_{k-1}$.
After finitely many steps, $F_N$ is affine (degree 1) and thus invertible.

\textbf{Step 3.} Then $F = F_N \circ \tau_N \circ \cdots \circ \tau_1$ is a composition of
invertible maps, hence invertible.

\medskip
The degree divisibility ensures that each peeling step is well-defined: if $\deg P > \deg Q$,
then $\deg P / \deg Q$ is an integer, and the leading form of $P$ is a power of the leading
form of $Q$ (by the Jacobian condition), allowing the peeling.
\end{proof}

\begin{remark}[Status of degree divisibility]
The degree divisibility is known for polynomial automorphisms of $\kk[x,y]$ (by
Jung--van der Kulk). Whether it holds for Keller maps not \emph{a priori} known to
be automorphisms is an \textbf{open question} whose resolution would settle JC$_2$.

The subleading term analysis (expand $\det J(P,Q) = c$ degree by degree after WLOG
normalizing $L = y$) shows that the Jacobian condition imposes algebraic relations
on the sub-leading coefficients but does \emph{not} immediately force $P$ and $Q$
to be univariate in $y$ (which would give the desired contradiction when $\deg P$
and $\deg Q$ are not in a divisibility relation). Explicit computation for
$(\deg P, \deg Q) = (5,3)$ shows that the $xy^3$ coefficient of $P_4$ and the $xy$
coefficient of $Q_2$ survive with the relation $a_0 = \frac{5}{3}b_0$.

Known degree constraints include $\gcd(\deg P, \deg Q) \geq 16$ (Lee--Li) and
$\gcd(\deg P, \deg Q) \neq 2p$ for any prime $p$, but these do not establish
divisibility.
\end{remark}

%% ====================================================================
\section{Connections to Lee--Li's Inner Polynomial Approach}
%% ====================================================================

Lee and Li \cite{LeeLi} study the 2-dimensional Jacobian Conjecture using Magnus' formula
for the formal inverse. Their key concept:

\begin{definition}[Inner polynomials, Lee--Li]
For a Keller map $F = (f, g)$, the \emph{inner polynomials} are the homogeneous components
$G_m$ of the formal inverse $G = F^{-1}$ (in $\kk[[x,y]]^2$), expressed via Magnus' formula
\[
  G = \sum_{m \geq 1} (-1)^{m-1} \frac{1}{m} \left(\frac{J(F)^{-1} - I}{1}\right)^{m-1}
    \cdot (f, g)^{\langle m \rangle}
\]
(notation from \cite{LeeLi}).
\end{definition}

\begin{theorem}[Lee--Li, Conjecture E]
Let $F = (f, g)$ be a Keller map with $\deg f = m > \deg g = n$, $m > n \geq 2$. If for
every inner polynomial $h_k$ (homogeneous component of $G$ of degree $k$), the northeastern
vertex of $N(h_k)$ lies in the region
\[
  \{(a, b) : a \leq k, \, b \leq k, \, a + b \leq k + 1\},
\]
then $F$ is a polynomial automorphism.
\end{theorem}

\begin{remark}
Lee--Li's Conjecture E, if proven, implies JC$_2$. The connection to our approach:
\begin{itemize}
\item Lee--Li's inner polynomials are the coefficients of the formal inverse $G$.
\item Our Lemma~\ref{lem:invlim} says: if $\deg G_{\leq d}$ stabilizes, then $G$ is
  polynomial and $F$ is an automorphism.
\item Lee--Li's Newton polygon constraints on $G_k$ are degree-bounding conditions that
  imply the stabilization in our sense.
\item The two approaches are therefore complementary: ours via the ``top-down'' peeling of
  Jung--van der Kulk decompositions, theirs via the ``bottom-up'' Magnus formula expansion.
\end{itemize}
\end{remark}

%% ====================================================================
\section{Concrete Research Questions}
%% ====================================================================

\begin{enumerate}
\item \textbf{Uniform degree bound:} For a Keller map $F = (P,Q)$ of degree $m$ in
  dimension 2, prove that the Anick approximants $\varphi_d$ can be chosen with
  $\deg \varphi_d \leq m$.

\item \textbf{Polydegree stabilization:} Prove Conjecture~\ref{conj:pdeg-stable}: the
  polydegree of the Anick approximants stabilizes for $d \gg 0$, implying that $F$
  is tame.

\item \textbf{Degree divisibility for Keller maps:} Prove that for a Keller map $(P,Q)$
  in dimension 2, either $\deg P \mid \deg Q$ or $\deg Q \mid \deg P$. This would
  reduce JC$_2$ to Theorem~\ref{thm:conditional}.

\item \textbf{Verify Zheglov's DC$_1$ proof:} Carefully verify the proof of DC$_1$ in
  \cite{Zheglov2024}. If correct, this settles JC$_2$ via the Belov-Kanel--Kontsevich
  equivalence.

\item \textbf{Lee--Li Conjecture E:} Prove the Newton polygon constraints on inner
  polynomials. This may be approachable via the root-subgroup structure of $\GJac(A_d)$.

\item \textbf{Properness from root structure:} Prove that a polynomial map expressible as an
  inverse limit of outer root products is necessarily proper.

\item \textbf{Connection to the counterexample:} Analyze how the tangent sweep
  mechanism in dimension $\geq 3$ exploits the existence of wild automorphisms
  (absent in dimension 2 by Jung--van der Kulk) to construct non-injective Keller maps.
  Note that inner roots of $\pp^d$ vanish in all dimensions; the obstruction in
  dimension $\geq 3$ is the failure of the tame = all identity, not root structure.
\end{enumerate}

%% ====================================================================
\section{Appendix: Summary of Notation}
%% ====================================================================

\begin{center}
\begin{tabular}{ll}
\hline
$A_d$ & $\kk[x_1,\ldots,x_n]/\pp^d$ \\
$\pp$ & maximal ideal $(x_1,\ldots,x_n)$ \\
$R(I)$ & Demazure roots of the monomial ideal $I$ \\
$R_{\mathrm{out}}, R_{\mathrm{inn}}$ & outer and inner roots \\
$U_\alpha$ & root subgroup $\cong \Ga$ \\
$\GJac(A_d)$ & $\{\varphi \in \Aut(A_d) : \det J(\varphi) \in \kk^\times\}$ \\
$T$ & maximal torus $\cong \Gm^n$ \\
$\divop$ & divergence operator \\
$J(F)$ & Jacobian matrix of $F$ \\
$N(P)$ & Newton polygon of $P$ \\
\hline
\end{tabular}
\end{center}

\begin{thebibliography}{99}
\bibitem{Alpoge2026}
L.~Alp\"oge, \emph{A counterexample to the Jacobian conjecture in dimension three},
preprint, July 2026.

\bibitem{Gallagher2026}
S.~Gao (following Gallagher), \emph{Counterexamples to the Jacobian conjecture in
dimensions greater than two}, arXiv:2608.00222, 2026.

\bibitem{DiazDP}
R.~D\'iaz, A.~Dubouloz, C.~Petitjean, \emph{Two variable monomial algebras with linear
automorphism groups}, preprint.

\bibitem{DiazDL}
R.~D\'iaz, A.~Dubouloz, A.~Liendo, \emph{Topologically integrable derivations and additive
group actions on affine ind-schemes}, preprint.

\bibitem{DiazLM}
R.~D\'iaz, G.~Lucchini Arteche, G.~Manzano-Flores, \emph{The structure of automorphism
groups of zero-dimensional monomial algebras}, preprint.

\bibitem{Anick}
D.~Anick, \emph{Limits of tame automorphisms of $k[x_1,\ldots,x_n]$}, J.~Algebra
\textbf{82} (1983), 459--468.

\bibitem{JungVdK}
H.~Jung, \emph{\"Uber ganze birationale Transformationen der Ebene}, J.~Reine Angew.\ Math.\
\textbf{184} (1942), 161--174.

\bibitem{AbhyankarMoh}
S.~Abhyankar, T.T.~Moh, \emph{Embeddings of the line in the plane}, J.~Reine Angew.\ Math.\
\textbf{276} (1975), 148--166.

\bibitem{Rentschler}
R.~Rentschler, \emph{Op\'erations du groupe additif sur le plan affine}, C.R.\ Acad.\ Sci.\
Paris S\'er.\ A \textbf{267} (1968), 384--387.

\bibitem{Moh}
T.T.~Moh, \emph{On the Jacobian conjecture and the configurations of roots}, J.~Reine
Angew.\ Math.\ \textbf{340} (1983), 140--212.

\bibitem{Nguyen}
T.~Nguyen, \emph{Some classes satisfying the 2-dimensional Jacobian conjecture and a proof
of the complex conjecture until degree 104}, arXiv:1902.05923, 2019.

\bibitem{LeeLi}
K.~Lee, L.~Li, \emph{On the two-dimensional Jacobian conjecture: Magnus' formula
revisited, IV}, arXiv:2408.01279, 2024.

\bibitem{BCW}
H.~Bass, E.~Connell, D.~Wright, \emph{The Jacobian conjecture: reduction of degree and
formal expansion of the inverse}, Bull.\ Amer.\ Math.\ Soc.\ (N.S.) \textbf{7} (1982),
287--330.

\bibitem{BKK}
A.~Belov-Kanel, M.~Kontsevich, \emph{The Jacobian conjecture is stably equivalent to the
Dixmier conjecture}, Mosc.\ Math.\ J.\ \textbf{7} (2007), no.~2, 209--218.

\bibitem{Tsuchimoto}
Y.~Tsuchimura, \emph{Endomorphisms of Weyl algebra and $p$-curvatures},
Osaka J.~Math.\ \textbf{42} (2005), 435--452.

\bibitem{Zheglov2024}
A.~Zheglov, \emph{Proof of the Dixmier conjecture for the first Weyl algebra},
arXiv:2410.06959, 2024 (v5, January 2026, 78 pages).

\bibitem{Furter}
J.-P.~Furter, \emph{On the length of polynomial automorphisms of the affine plane},
Math.\ Ann.\ \textbf{322} (2002), 401--411.

\bibitem{FurterKaras}
J.-P.~Furter, M.~Karas, \emph{On relations of leading terms},
arXiv:0808.1821, 2008.

\bibitem{MakarLimanov}
L.~Makar-Limanov, \emph{A new proof of the Jung theorem},
arXiv:1409.1872, 2014.

\bibitem{Tao2026}
T.~Tao, \emph{A digestion of the Jacobian conjecture counterexample}, blog post, July 2026.
\end{thebibliography}

\end{document}
